\documentclass[11pt]{article} 
\usepackage{fullpage}
\usepackage[left=1in,top=1in,right=1in,bottom=1in,headheight=3ex,headsep=3ex]{geometry}

\newcommand{\blankline}{\quad\pagebreak[2]}

\title{NCCU IMES \\ Causal Inference and Data Science in Economics}
\author{Tzu-Ting Yang and Po-Chun Huang}
\date{Spring 2026}

\usepackage[sc]{mathpazo}
\linespread{1.05}
\usepackage[T1]{fontenc}

\usepackage{setspace}
\usepackage{multicol}
\usepackage{fancyhdr,lastpage}
\usepackage{url}
\pagestyle{fancy}
\usepackage{hyperref}
\usepackage{amsmath}

\rhead{\footnotesize Causal Inference and Data Science in Economics -- Spring 2026}
\cfoot{\small \thepage/\pageref*{LastPage}}

\hypersetup{colorlinks,breaklinks,
            linkcolor=blue,urlcolor=blue,
            anchorcolor=blue,citecolor=black}

\begin{document}

\maketitle

\blankline

\begin{tabular*}{1.03\textwidth}{@{\extracolsep{\fill}}lr}

Instructor 1: Prof. Tzu-Ting Yang & Instructor 2: Prof. Po-Chun Huang \\

E-mail: \texttt{tty@g.nccu.edu.tw} & \texttt{huangpo5@nccu.edu.tw} \\

\multicolumn{2}{@{}l}{
\href{https://sites.google.com/view/cpelab/teaching/2026-causal-inference-and-data-science-in-economics}{Course Website}
} \\
Class Hours: Thursday. 16:10-19:00 & Class Room: 270206 General Purpose Building \\

\hline
\end{tabular*}

\vspace{10 mm}

%------------------------------------------------
\section*{Course Description}

This course will survey empirical methods for conducting causal inference and data science in economics. We focus on recent methodological developments and their empirical applications. Topics include randomized experiments, matching methods, regression and causal machine learning, differences-in-differences design, synthetic control methods, regression discontinuity design, instrumental variables, shift-share design, and geographic data analysis.

We emphasize practical implementation of these methods through writing a term paper. Students are required to use \textit{Latex} to type their term paper. The course also introduces the role of generative AI tools in data processing and empirical research workflow. After taking this course, students should be able to conduct empirical research independently and critically evaluate applied economics studies.

%------------------------------------------------
\section*{Suggested Readings}

\begin{itemize}
\item Huntington-Klein (2021), \textit{The Effect: An Introduction to Research Design and Causality}
\item Cunningham (2020), \textit{Causal Inference: The Mixtape}
\item Abadie and Cattaneo (2018), Econometric Methods for Program Evaluation
\item Angrist and Pischke (2016), \textit{Mastering Metrics}
\item Angrist and Pischke (2009), \textit{Mostly Harmless Econometrics}
\item James et al. (2017), \textit{An Introduction to Statistical Learning}
\item Korinek (2023), Generative AI for Economic Research
\item Korinek (2025), AI Agents for Economic Research
\end{itemize}

%------------------------------------------------
\section*{Statistical Softwares}

\begin{itemize}
\item STATA
\item R
\item Generative AI tools (for coding assistance and workflow design)
\end{itemize}

%------------------------------------------------
\section*{Grading Policy}

\begin{itemize}
\item Two empirical homework \& Two Oral Exams (30\%)
\item Reading group presentation (15\%)
\item Term paper presentation (15\%)
\item Term paper (40\%): milestones throughout the term
%\item Extra 5 bonus points if you upload replication files to \textit{GitHub}
\end{itemize}

%Each homework is followed by an individual oral exam (20--30 minutes). The oral exam evaluates your understanding of your own code, empirical design, regression specification, and identification strategy. After each oral exam, we will briefly discuss your term paper idea and refine your research design.

%------------------------------------------------
\section*{Important Dates}

\begin{itemize}
\item Homework 1: 4/12
\begin{itemize}
\item Oral Exam 1: 4/23 week
\end{itemize}
\item Homework 2: 5/17
\begin{itemize}
\item Oral Exam 2: 5/28 week
\end{itemize}
\item Reading Group Presentation: 5/28 and 6/4
\item Term paper presentation: 6/11
\item Term paper deadline: 6/18
\end{itemize}

%------------------------------------------------
\section*{Course Outline}

\begin{itemize}
\item Course Introduction
\item Selection Bias
\item Randomized Experiment 
\item Matching Methods
\item Regression and Causal Machine Learning
\item Fixed Effects Models
\item Differences-in-Differences Design
\item Synthetic Control Method
\item Regression Discontinuity Design
\item Instrumental Variables
\item Shift-Share IV Design
\item Spatial RD Design
\item Causal Forest
\item Data Cleaning and Research Workflow
\item Generative AI in Empirical Research
\end{itemize}

\end{document}